\section{Fun Fact}

\begin{frame}{\btitle{pA sample}}
    pA 的範測二跟範測三分別是 2026 TOI 一模跟二模成績

    然後二模第 $20$ 小剛好是 $67$，笑死

    怎麼一堆拿 $67$ 啊
\end{frame}

\begin{frame}{\btitle{pC}}
    \begin{itemize}
        \item {$12$ 個子任務，超級多}
    \end{itemize}
\end{frame}

\begin{frame}{\btitle{pD}}
    \begin{itemize}
        \item {模數給 1145141 是因為 114514\\然後還以為會用到質數的性質。}
        \item {OEIS A027423}
    \end{itemize}
\end{frame}

\begin{frame}{\btitle{pE}}
    \begin{itemize}
        \item {分蛋糕系列第四題}
        \item {什麼叫做複賽出來不會區間 DP 南區賽被打爆，所以pE出了區間dp的子題 \\ 那個人今年還是不會寫區間dp}

    \end{itemize}
\end{frame}

\begin{frame}{\btitle{Appendix}}
    因為不知道要放什麼在附錄

    所以拿會考數學參考公式放在上面。
\end{frame}

\begin{frame}{\btitle{賽中}}
    \item {Judge 太慢導致 pF 被卡}
    \item {pF 沒拿根本虧爛}
    \item {pD 有人直接把質數打在程式裡面}
    \item {pD 為什麼不用 bitset}
    \item {一堆 487}
\end{frame}
